\chapter{ニュートン宇宙論とフリードマン方程式}

前回は空間の計量について話した。次の問題は力学である。スケール因子は時間とともにどのように変化するのか。まず少し幾何学から始め、それからニュートン宇宙論に進む。フリードマン方程式はまだ一般相対性理論から導出されない。それはニュートンの方程式から導出される。

\section{一様性、等方性、そして時間依存の計量}

一様性と等方性は前提として置かれている。宇宙が一様かつ等方であり、時間とともに変化するならば、ハッブルの法則は常に成り立つ。任意の固定された時刻において、後退速度は距離に比例する。
\begin{equation}
v = H(t)\,d.
\end{equation}
比例関係は線形である。一定なのは時間ではなく空間である。同じ比例係数が、与えられた瞬間のどこでも当てはまる。

\begin{classroomqa}
\audiencequestion{それは線形に比例するという意味ですか。}\lecturetimestamp{00:02:06}
\lecturerresponse{はい。固定された時刻において、後退速度は距離に線形に比例する。比例係数はハッブル定数である。一般にはそれは時間に依存するが、空間の点ごとに変わることはない。}
\end{classroomqa}

だからこそ「``ハッブル定数''」という表現は危険である。それは空間内の位置に対して一定である。時間に対して一定ではない。

次に自然な疑問が生じる。空間が膨張するなら、なぜ物差しも膨張しないのか。答えは、物差しは原子スケールにおける宇宙論的膨張に伴ういかなるものよりもはるかに強い力によって結合されているということである。束縛された物体は束縛されたままでいる。膨張するのは、十分に切り離されたものの間の距離である。例えば、互いを強く結びつけておくには離れすぎている銀河のようなものである。

\begin{classroomqa}
\audiencequestion{測定棒も膨張しているという事実について何か問題がありますか。}\lecturetimestamp{00:04:00}
\lecturerresponse{いいえ。測定棒は膨張していない。座標が膨張している。物差しを結合している力は、膨張による効果よりもはるかに強いので、通常の束縛された物体は背景とともに膨張しない。十分に離れた銀河は確かに離れていくが、物差しは膨張しない。}
\end{classroomqa}

原子間に微小な余分の斥力効果を想像することもできる。それでも棒を引き裂くことはない。それは、小さな余分の引力が硬いばねの平衡位置を変えるのと同じように、平衡間隔をほんの少しずらすだけである。

\begin{classroomqa}
\audiencequestion{しかしそれを点の間の力としてモデル化することはできないのですね。そうでなければ加速度になってしまうから。}\lecturetimestamp{00:05:29}
\lecturerresponse{原子間の小さな余分の力は想像できるが、それが結合力に比べて微小であれば、平衡位置をわずかにずらすだけである。膨張それ自体は通常の力の法則を生み出さない。ある種の膨張は力の観点でモデル化できる。すなわち加速膨張であるが、それは後の問題である。}
\end{classroomqa}

したがって有用な座標は、膨張する物質分布に埋め込まれた座標である。余分な固有運動を持たない銀河は同じ座標値を保つ。時間とともに変化するのは、座標上の隔たりとメートルで測った物理的距離との関係である。その関係はスケール因子\(a(t)\)によって定式化される。

平坦な空間に対して、空間計量は次のように書ける。
\begin{equation}
ds^2 = a(t)^2\left(dx^2+dy^2+dz^2\right).
\end{equation}
ここで\(a\)は、座標距離の1単位だけ離れた隣接する点の間の、メートルで測った実際の距離である。宇宙が膨張すれば、\(a\)は時間依存になる。

時空へ移行するには、特殊相対論的な形
\begin{equation}
d\tau^2 = dt^2 - ds^2,
\end{equation}
から出発し、空間計量を代入する。これにより
\begin{equation}
d\tau^2 = dt^2 - a(t)^2\left(dx^2+dy^2+dz^2\right).
\end{equation}
が得られる。係数は空間内の位置に依存しない。それらは時間のみに依存する。これが計量の言葉での一様性の表明である。対角計量は次の形をとる。
\begin{equation}
g_{\mu\nu}=\mathrm{diag}\!\left(1,-a(t)^2,-a(t)^2,-a(t)^2\right).
\end{equation}
\editorialnote{ここに書かれた完全な対角行列は、00:14:33の周辺で語られた計量テンソルの記述に従っている。黒板の枠は線素を明確に示しているが、完全に判読可能な行列ではない。}

非対角項はない。\(dx\,dy\)もなく、\(dx\,dt\)もなく、他の混合項もない。計量は平坦な時空のものとほとんど同じくらい単純である。特異な点は、空間的な測度と時間との関係自体が時間依存であるということである。

\begin{figure}
\centering
\includegraphics[width=0.62\textwidth]{lecture_02_figure_02.png}
\caption{共動格子のスケッチと、時間依存のスケール因子を持つ時空間隔。}
\end{figure}

この描像が要点を捉えている。アリスとボブは共動座標では1単位だけ離れたままであるが、その1座標単位が表すメートル数は時間とともに増えていく。もし過去のある時点で\(a(t)\)がゼロであったなら、すべての物理的間隔がゼロになってしまう。あらゆるものが他のあらゆるものの上に押し潰される。だからこそ\(a\to 0\)は無限の密度に対応するのである。\(a\)が増大すれば、宇宙はますます希薄になる。

この計量の形はすべての座標系で有効なわけではない。銀河とともに動く座標系における形である。

\begin{classroomqa}
\audiencequestion{その方程式において、小さい\(t\)、\(a(t)\)、その\(t\)とは何か。}\lecturetimestamp{00:19:19}
\lecturerresponse{それは銀河とともに動く時計が測る時間である。同じ宇宙時計工場で作られ、銀河によって運ばれた同一の時計を想像してほしい。この座標系では、時間変数はまさにそれらの時計が測るものである。銀河の世界線に沿って、固有時間はこの座標時間と同じである。}
\end{classroomqa}

実際、共動銀河の世界線に沿って、
\begin{equation}
dx=dy=dz=0,
\end{equation}
したがって
\begin{equation}
d\tau^2 = dt^2,
\end{equation}
よって
\begin{equation}
dt=d\tau.
\end{equation}

\begin{classroomqa}
\audiencequestion{この形の計量を不変に保つ変換はあるか。}\lecturetimestamp{00:23:08}
\lecturerresponse{回転はそうである。なぜなら\(dx^2+dy^2+dz^2\)は回転の下で不変だからである。しかしローレンツブーストはこの形を保たない。物質そのものが銀河とともに動く座標系を選び出しており、他の座標系では世界は等方的に見えないであろう。}
\end{classroomqa}

\section{重力場中の投射体}

ニュートンを宇宙全体に適用する前に、より単純な問題に立ち返ると役に立つ。重い質量\(M\)、はるかに小さい質量\(m\)、そして外向きの初速の力積である。力積の後、小さい質量は単に重力場の中を運動し、ニュートンがそうなるはずだと言うとおりに振る舞う。

宇宙論の議論にとって重要なのは、軌道問題の全体ではない。エネルギー保存である。運動エネルギーは
\begin{equation}
\frac{1}{2}mv^2.
\end{equation}

\begin{figure}
\centering
\includegraphics[width=0.42\textwidth]{lecture_02_figure_03.png}
\caption{外向きの試験質量と運動エネルギーの項\(\frac{1}{2}mv^2\)。}
\end{figure}

重力ポテンシャルエネルギーは、二つの質量の積、ニュートンの定数、そして分母に距離の1乗を含む。板書では因子が\(mMG/D\)の順に現れるので、ここでもその順序を保つのがよい。
\begin{equation}
-\frac{mMG}{D}.
\end{equation}
したがって全エネルギーは
\begin{equation}
\frac{1}{2}mv^2-\frac{mMG}{D}.
\end{equation}
エネルギー保存が言うのは、問題としている特定の粒子についてこの量が時間とともに変化しないということだけである。異なる粒子は異なるエネルギーを持ちうる。

\begin{figure}
\centering
\includegraphics[width=0.62\textwidth]{lecture_02_figure_04.png}
\caption{動径方向に運動する試験質量の全ニュートンエネルギー：運動項から重力ポテンシャルを引いたもの。}
\end{figure}

負の符号が重要である。重力ポテンシャルエネルギーが負なのは、重力が引力だからである。粒子を近づけると、より大きな負のポテンシャルエネルギーが生じる。遠ざけるとポテンシャルエネルギーは負でなくなる。

\begin{classroomqa}
\audiencequestion{ポテンシャルエネルギーが距離とともに減少するのではなく増加するのか。}\lecturetimestamp{00:34:13}
\lecturerresponse{そのとおり。定数を忘れて\(1/D\)と\(-1/D\)のグラフを見よ。力はポテンシャルエネルギーが減少する向きに押す。したがって引力は逆さまの曲線\(-1/D\)に対応し、ポテンシャルエネルギーは距離とともに増加し、粒子が内側へ動くにつれてより負になる。}
\end{classroomqa}

\begin{figure}
\centering
\includegraphics[width=0.46\textwidth]{lecture_02_figure_05.png}
\caption{重力ポテンシャルエネルギーの符号に関する\(1/D\)と\(-1/D\)の定性的比較。}
\end{figure}

\section{脱出速度}

さて、可能性を分類しよう。エネルギーは正、負、またはゼロでありうる。この三つの場合しかない。

\begin{itemize}
\item エネルギーが負なら、粒子は束縛されている。小さすぎる速度で投げ上げれば、それは上昇し、減速し、折り返し、落ちてくる。
\item エネルギーが正なら、粒子は脱出する。重力は最初はそれを減速するが、止めるには足りず、やがてゼロでない漸近的な外向き速度で遠ざかっていく。
\item エネルギーがちょうどゼロなら、粒子は脱出速度にある。それが脱出するとしないとの間の境界の場合である。
\end{itemize}

初速度が十分小さければ、ポテンシャル項が最初から支配的であり、エネルギーは負である。

\begin{figure}
\centering
\includegraphics[width=0.62\textwidth]{lecture_02_figure_06.png}
\caption{小さな外向き速度では、負のポテンシャル項が全エネルギーを支配する。}
\end{figure}

ゼロエネルギーの場合、距離は増え続けるが、速度は減り続ける。粒子は有限の時間で静止することはできない。もし静止すれば、それは落ちて戻ってしまうであろう。したがって距離は非有界である一方、速度はゼロに近づいていく。

これが宇宙論にとって重要な直接的な類推である。宇宙も同じように振る舞うことができる。初期の外向きのハッブル流れが十分に大きければ、すべては他のすべてから逃げ出す。それが小さすぎれば、重力が勝ち、巨大な収縮が起こる。ちょうど脱出速度に完全に対応する境界的な場合もある。

\section{ニュートンの定理}

宇宙論的な簡約がうまくいくのは、逆二乗則に対するニュートンの定理のおかげである。

物質がある中心のまわりに球対称に分布しているなら、その分布の外側にある粒子に働く重力は、すべての質量が中心に集中している場合とまったく同じである。しかしこの定理はそれよりも強い。粒子が分布の内側にあるなら、より小さい半径にある質量だけが寄与する。より大きい半径にある質量は正味の力を及ぼさない。

この定理の最もすっきりした形は薄い殻である。完全な球殻は、その内側のどこにある粒子にも、中心から離れた位置であっても、重力を及ぼさない。これは一見すると驚くべきことだが、\(1/r^2\)則に対しては正しい。

\begin{classroomqa}
\audiencequestion{そちらの例で、粒子が境界を横切って動くにつれて力が不連続になるように見えるのは、少し驚くべきことではないでしょうか。}\lecturetimestamp{00:48:19}
\lecturerresponse{はい。理想的な無限に薄い殻に対しては驚くべきことであり、それは\(1/r^2\)力の法則に特有のことである。現実の殻には厚みがあるので、変化は滑らかになるが、理想的な殻についての鋭い主張は、まさにニュートンの定理の注目すべき性質である。}
\end{classroomqa}

そこで図を描くと定理がもっともらしく見えてくる。点質量を水の点源に置き換える。三次元では外向きの流れは\(1/r^2\)のように減少する。そのような源を球殻のまわりに連続的に置く。殻の外側では、流れはすべてが中心に集中している場合とまったく同じに見える。殻の内側では正味の流れはありえない。中心に源がないからである。同じ構造が重力の定理の基礎になっている。

ここでの要点は定理を再導出することではない。それを用いることである。

\section{一体系問題としての一様な宇宙}

さて宇宙論に戻ろう。銀河の一様な気体で満たされた世界の小さな領域を取る。望むなら自分を中心に置いてもよい。それは単なる便宜にすぎない。誰でも自分を中心と呼ぶことができ、同じ方程式が出てこなければならない。

\(x\)軸上で共動座標\(x\)にある銀河を選ぶ。宇宙は等方であるから、軸の選び方は問題にならない。\(x=0\)を中心とし、その銀河を通る球を描く。ニュートンの定理により、その銀河に働く力は、その球の内側のすべての質量が中心に集中している場合とまったく同じであり、外側のすべては何も寄与しない。

球は銀河とともに膨張する。その物理的半径は時間とともに変化するが、球は常に同じ銀河を含み、したがって同じ全質量を含む。

したがってその銀河の運動は、発射体とまったく同じ種類のエネルギー方程式に従う：
\begin{equation}
\frac{1}{2}mv^2-\frac{mMG}{D}=\text{constant}.
\end{equation}
ここで\(M\)は、ウォームアップの図における中心の質量ではなく、膨張する球の内側の質量を意味する。次に\(2\)を掛けて\(m\)で割ることで代数が簡単になる。定数を\(m\)で割ったものもやはり定数であるから、
\begin{equation}
v^2-\frac{2MG}{D}=\text{constant}.
\end{equation}
と書いてよい。ここでも``定数''とは、その特定の銀河について時間に対して一定であることを意味する。

次にこれをスケール因子と関係づける。銀河までの物理的距離は
\begin{equation}
D=ax,
\end{equation}
である。なぜなら\(x\)は固定された共動座標だからである。その速度は
\begin{equation}
v=\dot D=\dot a\,x.
\end{equation}

球の内側の質量は密度と体積の積である。\(\rho\)を質量密度とすると、
\begin{equation}
M=\frac{4\pi}{3}\rho D^3=\frac{4\pi}{3}\rho a^3x^3.
\end{equation}
エネルギー方程式に代入すると
\begin{equation}
\dot a^2x^2-\frac{2G}{ax}\left(\frac{4\pi}{3}\rho a^3x^3\right)=\text{constant},
\end{equation}
または
\begin{equation}
\dot a^2x^2-\frac{8\pi G}{3}\rho a^2x^2=\text{constant}.
\end{equation}

左辺の両方の項は\(x^2\)に比例する。そこにある他のすべてはどこでも同じである。したがって右辺も\(x^2\)に比例しなければならない。それはそれを定数と呼ぶことと矛盾しない。それは各固定された\(x\)について時間に対して一定であることを意味するが、銀河ごとに\(x^2\)とともにスケールしてもよい。それは理にかなっている。より遠い銀河はより速く動いている。

そこで共通因子\(x^2\)を割って除く。残った定数はもはや時間にも空間にも依存しない。それを歴史的に標準的な仕方で\(-k\)と書く：
\begin{equation}
\dot a^2-\frac{8\pi G}{3}\rho a^2=-k.
\end{equation}
\(a^2\)で割った後、
\begin{equation}
\left(\frac{\dot a}{a}\right)^2=\frac{8\pi G}{3}\rho-\frac{k}{a^2}.
\end{equation}

これは有名なフリードマン・ロバートソン・ウォーカー方程式であるが、講義ではこれが実はフリードマンの方程式であるとも述べられている。符号の規約は歴史的なものである。\(-k\)と書くことは\(k\)が正でなければならないことを意味しない。\(k\)は正、負、またはゼロでありうる。光速度を\(1\)に等しく設定する単位系では、\(k\)は無次元である。

\section{密度と物質優勢方程式}

次の疑問はすぐに浮かぶ。\(\rho\)とは何か。それは定数ではない。宇宙が膨張するにつれて、密度は減少する。

その依存性を表すために、座標空間に単位立方体を取る。\(x\)、\(y\)、\(z\)の各方向に一辺が1単位である。立方体は座標とともに膨張するので、同じ銀河は常にその立方体の内部にとどまる。したがって、その座標立方体の内部の質量は固定されている。

\editorialnote{講義では記号\(M\)をいくつかの段階で再利用している。ここでは\(M_{\Box}\)は、単位座標立方体内の質量を、以前の中心質量および球に囲まれた質量と区別するためにのみ用いられている。}

立方体の物理的な一辺の長さは\(a\)であるから、その物理的な体積は\(a^3\)である。したがって
\begin{equation}
\rho=\frac{M_{\Box}}{a^3}.
\end{equation}
これは物質優勢の場合である。右辺には通常の非相対論的物質の密度のみが現れる。

\begin{classroomqa}
\audiencequestion{その方程式によれば、ハッブル定数は\(\rho\)の関数である。したがって、それは定数ではない。}\lecturetimestamp{01:19:21}
\lecturerresponse{そのとおり。ハッブル定数は定数ではない。それは密度に依存し、密度は時間とともに変化する。講義では、期末試験が一問だけ、すなわち「ハッブル定数は定数か」という問題になりうると冗談まで言っている。}
\end{classroomqa}

\(\rho=M_{\Box}/a^3\)をフリードマンの方程式に代入すると
\begin{equation}
\left(\frac{\dot a}{a}\right)^2=\frac{8\pi G M_{\Box}}{3a^3}-\frac{k}{a^2}.
\end{equation}

この時点で、任意の選択肢が明らかになる。単位座標立方体内の質量と\(k\)の値が必要である。\(k\)の取りうる三つの符号は、投射体問題の三つのエネルギー場合に対応する。

\begin{itemize}
\item もし\(-k>0\)ならば、エネルギーは実質的に脱出速度を上回っており、宇宙は永遠に膨張する。
\item もし\(-k<0\)ならば、重力が最終的に勝ち、宇宙は再収縮する。
\item もし\(k=0\)ならば、宇宙はちょうど脱出速度にある。
\end{itemize}

\begin{classroomqa}
\audiencequestion{\(k=0\)か。すると二つの銀河の間の速度はゼロに近づくのか。}\lecturetimestamp{01:21:12}
\lecturerresponse{そのとおり。臨界の場合、速度はゼロに近づくが、距離は依然として非有界である。銀河は分離し続けるが、ただますますゆっくりとである。もし有限の時間で実際に静止するならば、それらは逆戻りして落ち込むであろう。}
\end{classroomqa}

\section{臨界の場合\(k=0\)}

ここで最も単純な場合、\(k=0\)を選ぶ。これはちょうど脱出速度の場合であり、数学が単純であるという理由だけで単純なのである。方程式は
\begin{equation}
\left(\frac{\dot a}{a}\right)^2=\frac{8\pi G M_{\Box}}{3a^3}.
\end{equation}
となる。

右辺は常に正であり、決してゼロにならない。したがって\(\dot a/a\)はゼロを通過できない。この場合、宇宙は向きを変えられない。常に膨張しているか、常に収縮しているかのどちらかである。膨張する枝を取る。

ここで冪法則を推測する。
\begin{equation}
a(t)=ct^p.
\end{equation}
すると
\begin{equation}
\dot a = cp\,t^{p-1},
\end{equation}
したがって
\begin{equation}
\frac{\dot a}{a}=\frac{p}{t}.
\end{equation}
微分方程式に代入すると
\begin{equation}
\frac{p^2}{t^2}=\frac{8\pi G M_{\Box}}{3c^3t^{3p}}.
\end{equation}
\(t\)の冪が一致するためには
\begin{equation}
3p=2.
\end{equation}
でなければならない。よって
\begin{equation}
p=\frac{2}{3}.
\end{equation}
次に定数は
\begin{equation}
\frac{4}{9}=\frac{8\pi G M_{\Box}}{3c^3},
\end{equation}
を満たさなければならず、ゆえに
\begin{equation}
c^3=6\pi G M_{\Box}.
\end{equation}
したがって
\begin{equation}
a(t)=\left(6\pi G M_{\Box}\right)^{1/3}t^{2/3}.
\end{equation}

依然として時間の加法的なずれが存在しうる。それは\(t\)を\(t-t_0\)に置き換えることによって吸収される。\(t=0\)で\(a=0\)となるように時間の原点を選べば、その残りの定数は消える。

したがって宇宙は時間の3分の2乗として膨張する。
\begin{equation}
a(t)\propto t^{2/3}.
\end{equation}
時間に線形には膨張しない。線形の成長は、加速度のない運動の類似物であろう。ここでは膨張は減速する。

式を\(t=0\)まで遡って押し進めると、スケール因子はゼロになり、密度は発散する。\(a=0\)と言うのは極端である。物質があまりに高密度になったときに何が起こるかはわかっていない。本当の要点は、ある初期の段階で、原子、原子核、クォークがすべて異常に高密度の状態に詰め込まれていたということである。

\section{ハッブル定数と宇宙の年齢}

臨界模型では
\begin{equation}
H(t)=\frac{\dot a}{a}=\frac{2}{3t}.
\end{equation}
したがってハッブル定数はほぼ宇宙の年齢の逆数である。
\begin{equation}
t=\frac{2}{3H}.
\end{equation}
これは計算の主要な成果の一つである。今日のハッブル定数を測定すれば、少なくともニュートン模型が意味をなさなくなる時点までは、宇宙の年齢の推定値が得られる。

引用されている年齢はおよそ\(13\)億年から\(14\)億年である。また、\(k=0\)は観測的に非常に現実的であるとも述べられている。そこに引用されている精度の水準では、\(k=0\)からの有意なずれは知られていない。

\begin{classroomqa}
\audiencequestion{私は\(M\)と\(a\)が何なのか、ただ少し混乱しているだけです。}\lecturetimestamp{01:39:16}
\lecturerresponse{どちらも格子の選び方による任意性を含んでいる。基準となる箱を変えれば、その中の物質の量も変わり、スケール因子の数値も変わる。物理的なのは密度\(M/a^3\)、あるいはそれと等価な、格子の選び方に依存しない\(a/M^{1/3}\)のような組み合わせである。\(k=0\)解では、\(t^{2/3}\)に掛かる未知の定数は\(\dot a/a\)から消えるので、年齢の関係式\(t=2/3H\)は基準となる箱の質量に依存しない。}
\end{classroomqa}

その相殺は解において直ちに示される。もし
\begin{equation}
a(t)=Ct^{2/3},
\end{equation}
ならば
\begin{equation}
\dot a=\frac{2}{3}Ct^{-1/3},
\end{equation}
であり、したがって
\begin{equation}
\frac{\dot a}{a}=\frac{2}{3t}.
\end{equation}
となる。定数\(C\)は完全に消える。

\section{ニュートン的な導出がうまくいく理由}

終わり近くで問題はこうなる。なぜニュートン的な導出は、これほどうまく相対論的な導出と一致するべきなのか。

\begin{classroomqa}
\audiencequestion{これはすべて、開いた宇宙という仮定に依存しているように思えます。}\lecturetimestamp{01:45:44}
\lecturerresponse{この段階では導出はニュートン的であり、ニュートン物理学は平坦な空間を意味する。\(k\)を空間の幾何学として解釈することは次の講義に先送りされる。そこで\(k=0\)は平坦な空間に対応し、他の符号は曲がった空間幾何学に結び付けられる。}
\end{classroomqa}

\begin{classroomqa}
\audiencequestion{ずっと昔、メートル尺の長さは\(a\)に依存していたのでしょうか。}\lecturetimestamp{01:47:53}
\lecturerresponse{十分に昔には、物質はあまりに密集していたため、別々のメートル尺は実際には存在しなかっただろう。それより前には、それらは押し合わされていたはずである。この講義はまた、相対論はそもそも真に剛体な物体を許さないことを思い起こさせる。もし物体が完全に剛体なら、片側を押すと信号が瞬時に、光より速く反対側へ伝わることになる。}
\end{classroomqa}

\begin{classroomqa}
\audiencequestion{なぜニュートン的と相対論的で、このように方程式が一致するのだと思いますか。}\lecturetimestamp{01:49:03}
\lecturerresponse{ニュートン的な扱いは、十分に小さな領域にのみ適用されているからである。その領域が宇宙全体の曲率半径に比べて小さく、そこでの局所的なハッブル速度が光速より十分小さい限り、ニュートン力学は良い近似である。この講義ではこれを、地球上で草の成長を研究することと比較している。地球全体の曲率を使わなくても、その局所的な過程は理解できる。}
\end{classroomqa}

これが局所領域の論理である。ニュートン力学は、第一に、空間の曲率が当面の問題にとって重要でないこと、第二に、局所的に光速に近い速度で動くものがないことを仮定する。宇宙論では、それは均質であるのに十分大きく、しかし曲率の効果が無視でき、ハッブル速度が非相対論的にとどまるほどにはまだ十分小さい領域を選ぶことを意味する。

そう考えると、ニュートン的計算は偶然ではない。それは、相対論的な形でも読める方程式の局所近似である。曲率の項を左辺に移せば、フリードマン方程式は
\begin{equation}
\left(\frac{\dot a}{a}\right)^2+\frac{k}{a^2}=\frac{8\pi G}{3}\rho.
\end{equation}
となる。左辺は幾何学的である。それはスケール因子と、次回には空間の曲率を含む。右辺はエネルギー密度である。これはまさに一般相対論から期待される種類の方程式である。

次回の主題は曲率そのもの、アインシュタイン方程式、そしてその相対論的導出とニュートン的導出との関係である。
